\section{Affine neighbourhoods for factors of the SNC total transform}

\begin{lemma}[One factor on an adapted affine neighbourhood]\label{lem:sc0-snc-affine-factor-neighbourhood}
\lean{EnvelopeProperty.Char0.exists_factor_stalk_monomial_in_component_sections,EnvelopeProperty.Char0.exists_snc_etale_chart_monomial_factor}
\leanok
\uses{lem:sc0-snc-factor-stalks,lem:sc0-snc-etale-charts}
Let $X$ be smooth of relative dimension $n$ over a perfect field, let $F$ be a finite simple
normal crossing family, and suppose $AB=\prod_iF_i^{e_i}$ for ideal sheaves $A,B$.
At every closed point $x$ there is an affine neighbourhood with \'etale coordinate functions
$v_1,\ldots,v_n$ such that the members of $F$ through $x$ are distinct coordinate hyperplanes,
all other members are absent, and $A$ is generated by a monomial in those same assigned
coordinate functions with exponents bounded by $e_i$.
\end{lemma}
\begin{proof}
Lemma~\ref{lem:sc0-snc-factor-stalks} gives regular parameters $z$, the injective component
assignment $c$, and bounded exponents $a_i$ with
$A_x=(\prod_{i:x\in F_i}z_{c(i)}^{a_i})$. The number of parameters equals $n$, since $x$ is closed.
Spread this regular system to affine \'etale coordinates adapted to $F$; let $c'$ be their
component assignment. At the stalk, $(z_{c(i)})=(F_i)_x=(v_{c'(i)})$. Taking powers and products
therefore preserves the monomial ideal. Thus $A_x$ is generated by the germ of the regular
function $M=\prod_{i:x\in F_i}v_{c'(i)}^{a_i}$ on this affine chart.
Its coordinate ring is Noetherian. Equality of the two finitely generated ideals after
localization at $x$ spreads to equality on a principal open neighbourhood of $x$. Restrict the
coordinate functions and all component identities there; restriction preserves the \'etale
coordinate map, absence of the other components, and the monomial product formula.
\end{proof}

\begin{lemma}[Simultaneous factors on an adapted affine neighbourhood]\label{lem:sc0-snc-simultaneous-neighbourhood}
\lean{EnvelopeProperty.Char0.exists_snc_etale_chart_monomial_family}
\leanok
\uses{lem:sc0-snc-factor-stalks,lem:sc0-snc-affine-factor-neighbourhood}
For finitely many ideal sheaves $A_j$ with $\prod_j A_j=\prod_i F_i^{e_i}$ as above, every closed
point has one adapted affine \'etale neighbourhood on which all $A_j$ are generated by monomials
in the same assigned boundary coordinates.
\end{lemma}
\begin{proof}
For each $j$, write the total product as $A_jB_j$, where $B_j$ is the product of the other
factors. All $A_j$ are invertible, since their product is invertible. Fix one regular system
of parameters and one component assignment for $F$ at the point. Each stalk generator divides
the total monomial, so prime factorization gives its bounded monomial exponents in this fixed
system. Spread the common SNC parameters to one adapted \'etale chart. The preceding comparison
of component ideals identifies every factor stalk with its corresponding coordinate monomial.
Spread these finitely many ideal equalities simultaneously by inverting the product of the
finitely many denominators. The empty family requires just the adapted \'etale chart.
\end{proof}

\begin{lemma}[A finite simultaneous monomial cover]\label{lem:sc0-snc-simultaneous-finite-cover}
\lean{EnvelopeProperty.Char0.exists_finite_snc_etale_monomial_cover}
\leanok
\uses{lem:sc0-snc-simultaneous-neighbourhood}
If $X$ is quasi-compact, the simultaneous adapted affine neighbourhoods above admit a finite
subcover of $X$.
\end{lemma}
\begin{proof}
The scheme is Jacobson, because it is locally of finite type over a field. Every point specializes
to a closed point, and therefore lies in the adapted open neighbourhood of that closed point.
These neighbourhoods cover $X$; quasi-compactness gives a finite subcover.
\end{proof}
